\documentclass[11pt]{article}

\usepackage[margin=1in]{geometry}
\usepackage{amsmath,amssymb,booktabs,tabularx}

\title{Open-news vs closed-gap decomposition}
\author{}
\date{}

\begin{document}
\maketitle

\noindent\textbf{Status.} Draft\\
\textbf{Plan item.} Open-news vs closed-gap decomposition\\
\textbf{Depends on.} \texttt{13\_news\_mark\_mgf.tex}\\
\textbf{Feeds.} News and gaps section, sufficient-statistics consistency, gap-behavior validation

\section{Purpose}

Derive the decomposition that routes the same coupled news return mark into an
open-market news jump \(M_t\) or a reopening gap \(G_t\), depending only on the
market state at the event timestamp.
This makes explicit why Janiform treats open news and closed-session news as one
mark law with different observation timing, not as two unrelated return-shock
mechanisms.

\section{Setup}

Let \(c\in\{\mathrm{known},\mathrm{surprise}\}\) index the news component, and
let \(\{\tau_m^c\}\) be its event times.
Each event has a shared severity mark \(U_m^c\) and a coupled return-jump
mark \(\zeta_m^c\).
The activity shock always enters the news activity accumulator,
\[
J_t^c=\sum_{\tau_m^c\le t}a_cU_m^c.
\]

Let \(O(\tau)\) be the open-market indicator at event time \(\tau\), and let
\[
C(\tau)=1-O(\tau)
\]
be the closed-market indicator.
For a closed-market event, let \(r(\tau)\) be the first reopening time
after \(\tau\).
For an open-market event, \(r(\tau)\) is not used.
Open-market news return jumps are applied through \(M_t\).
Closed-market news return jumps are stored until
reopening and applied through \(G_t\).

\section{Notation}

\begin{tabularx}{\textwidth}{lXX}
\toprule
Symbol & Meaning & Units or domain\\
\midrule
\(\tau_m^c\) & News event time for component \(c\) & Calendar time\\
\(O(\tau)\) & Open-market event indicator & \(\{0,1\}\)\\
\(C(\tau)\) & Closed-market event indicator & \(\{0,1\}\)\\
\(r(\tau)\) & First reopening time after a closed event & Calendar time\\
\(U_m^c\) & Shared severity mark & Non-negative\\
\(\zeta_m^c\) & Coupled return-jump mark & Real\\
\(J_t^c\) & Activity-shock accumulator & Non-negative\\
\(M_t^c\) & Open-market news return-jump accumulator & Real\\
\(G_t^c\) & Closed-market reopening-gap accumulator & Real\\
\bottomrule
\end{tabularx}

\section{Assumptions used}

\begin{tabularx}{\textwidth}{XXX}
\toprule
Assumption & Where it is used & Why it is needed\\
\midrule
Every news event has one market-state classification at its event timestamp.
& Splitting each event with \(O(\tau)\) and \(C(\tau)\).
& Prevents the same return mark from being counted as both open news and a closed gap.\\
Open and closed indicators are complementary.
& Showing disjointness of \(M_t\) and \(G_t\) routing.
& Ensures \(O(\tau)+C(\tau)=1\).\\
Closed-session return jumps are observable at first reopening.
& Defining \(G_t\) using \(r(\tau)\).
& Connects closed news to realized close-open gap returns.\\
The same coupled mark \(\zeta_m^c\) is used in either route.
& Writing \(M_t^c\) and \(G_t^c\).
& Keeps open-news jumps and closed gaps as different timing views of one news-return mechanism.\\
Activity shock routing is independent of open or closed return routing.
& Defining \(J_t^c\) without \(O(\tau)\) or \(C(\tau)\).
& Closed news can still excite later activity even though its return jump is delayed to reopening.\\
\bottomrule
\end{tabularx}

\section{Result}

\noindent\textbf{Identity.} For each news component \(c\),
\[
M_t^c
=
\sum_{\tau_m^c\le t}O(\tau_m^c)\zeta_m^c,
\qquad
G_t^c
=
\sum_{r(\tau_m^c)\le t}C(\tau_m^c)\zeta_m^c.
\]
The total open-news and closed-gap accumulators are
\[
M_t=M_t^{\mathrm{known}}+M_t^{\mathrm{surprise}},
\qquad
G_t=G_t^{\mathrm{known}}+G_t^{\mathrm{surprise}}.
\]
For one event at time \(\tau\), with \(r(\tau)\) used only when
\(C(\tau)=1\), the return mark contribution is routed as
\[
\Delta M_\tau^c=O(\tau)\zeta^c,
\qquad
\Delta G_{r(\tau)}^c=C(\tau)\zeta^c,
\qquad
\Delta M_\tau^c\Delta G_{r(\tau)}^c=0.
\]

\section{Derivation}

For a single event at time \(\tau\), the market is either open or closed, so
\[
O(\tau)\in\{0,1\},
\qquad
C(\tau)=1-O(\tau),
\qquad
O(\tau)C(\tau)=0.
\]
The same coupled return mark \(\zeta^c\) must be observed through exactly one
return channel.
Therefore, the open-channel increment at the event time is
\[
\Delta M_\tau^c=O(\tau)\zeta^c.
\]
If \(O(\tau)=1\), this equals \(\zeta^c\). If \(O(\tau)=0\), it equals zero.

For a closed event, the return mark is not an open-market direct jump at
\(\tau\).
It is carried to the first reopening time \(r(\tau)\), giving
\[
\Delta G_{r(\tau)}^c=C(\tau)\zeta^c.
\]
If \(C(\tau)=1\), this equals \(\zeta^c\) at reopening. If \(C(\tau)=0\), it
equals zero.

Summing open-channel increments over events whose event time is at or before
\(t\) gives
\[
M_t^c
=
\sum_{\tau_m^c\le t}O(\tau_m^c)\zeta_m^c.
\]
Summing closed-channel increments over events whose reopening time is at or
before \(t\) gives
\[
G_t^c
=
\sum_{r(\tau_m^c)\le t}C(\tau_m^c)\zeta_m^c.
\]
Since \(O(\tau)C(\tau)=0\), no single event can contribute a nonzero return
mark to both channels:
\[
\Delta M_\tau^c\Delta G_{r(\tau)}^c
=
O(\tau)C(\tau)(\zeta^c)^2
=0.
\]

The activity shock is not decomposed by return-jump timing.
It is accumulated from every event:
\[
J_t^c
=
\sum_{\tau_m^c\le t}a_cU_m^c.
\]
Thus, the news event has one shared severity channel for activity and one
return-jump mark routed by market state.

\section{Domain and validity}

The decomposition requires a deterministic market-state classifier for event
timestamps, complementary open and closed states, and a well-defined reopening
map for closed events.
It also requires finite return marks and non-negative severity marks.
The identity is a routing and accounting result, not a claim that closed-market
price discovery is directly observed before reopening.

\section{Consequences}

\begin{itemize}
    \item Open news and closed gaps share the same return-mark law and differ by market-state routing.
    \item Closed events do not create an \(M_t\) jump at the closed timestamp.
    \item Open events do not create a \(G_t\) gap contribution.
    \item Every active news return jump is paired with a news activity shock through the shared severity mark.
    \item The decomposition supports separate diagnostics for open jumps and closed gaps without introducing a second news-return mechanism.
\end{itemize}

\section{Checks}

\begin{tabularx}{\textwidth}{XXX}
\toprule
Check & Expected result & Evidence location\\
\midrule
Shared dispatch routing
& A closed event increments \(G_t\) and not \(M_t\); an open event increments \(M_t\) and not \(G_t\).
& Test \texttt{news\_dispatch\_routes\_same\_coupled\_return\_jump\_by\_market\_state} in \texttt{crates/janiform-model-kernels/tests/news\_coupling\_distribution\_regression.rs}\\
Canonical dispatch implementation
& Dispatch sorts events deterministically, increments \(J_t\) with the constructed activity-shock mark for every event, and routes the coupled return jump by market state.
& \texttt{dispatch\_shared\_news\_events} in \texttt{crates/janiform-model-kernels/src/news/dispatch.rs}\\
Fixture validation
& Open-market expected outcomes must have zero closed-gap increment, and closed-market expected outcomes must have zero open-jump increment.
& \texttt{\_validate\_open\_closed\_gap\_behavior\_expected\_outcome} in \texttt{tests/support/fixture\_validation.py}\\
Gap behavior diagnostics
& Controlled open/closed fixtures produce separate open-jump and closed-gap summaries with the expected totals.
& Test \texttt{gap\_behavior\_metrics\_match\_controlled\_open\_closed\_fixture} in \texttt{crates/janiform-validation/tests/gap\_behavior\_diagnostics.rs}\\
Regression harness
& Open-window closed-gap leakage is reported in the closed-gap semantic bucket.
& Test \texttt{test\_open\_closed\_gap\_leakage\_regression\_harness\_reports\_semantic\_bucket\_for\_failures} in \texttt{tests/test\_open\_closed\_gap\_leakage\_regression\_harness.py}\\
\bottomrule
\end{tabularx}

\end{document}
